% !TeX root = ../papers/nc-p3.tex
\section{重複度込みの十九次での二次曲面モデル}
\label{sec:sectional-nineteen-model}
\begin{theorem}\label{thm:sectional-nineteen-model}
定理\ref{thm:sectional-cofactor-length}で $d=19$ とする。核像は整二次曲面 $Q$ で、
核写像は双有理である。滑らかな場合は $S_0=Q$、錐の場合は最小解消
$S_0=\mathbb F_2$ とし、核超平面を $L$ と書く。ある母線束 $F$ に対し、
入力逆写像は $A=2L+F$ の四切断で定まる。共通基点スキームは長さ四の局所完全交叉である。
階数一軌跡には必ず曲線がある。従ってその被約曲線部分は直線、整二次曲線、
交わる二直線のいずれかである。
\end{theorem}
\begin{proof}
曲線核は定理\ref{thm:genus-two-kernel-excess}で除かれた。
$HL=5$, $H^2=8$ よりHodge指数から $2\le L^2\le3$、核は非退化なので
二次または三次曲面への双有理写像しかない。固定部を除いた逆系 $A$ は滑らかな最小モデルで
$AL=5$, $A^2=8+\sum m_i^2$ を満たす。ここで $m_i$ は無限近接を含む基点の重複度である。
また一般平面での核像 $J$ は自己準同型曲線 $G$ と同型で、算術種数二である。
この事実を使えば定理\ref{thm:cofactor-nineteen-surface}の三次曲面の数値排除がそのまま成立する。
詳しくは $A^2\le25/3$ より基点がなく、正規非錐では最小解消の $K=-L$ に対し
$A(A+K)=8-5$ の偶奇が矛盾する。錐の正規化母線束は $A=aE+5F$ で
$10a-3a^2=8$ が $a=2$ を強制し $AE<0$。非正規非錐では
$S_0=\mathbb F_1$, $A=2E+3F$ で一般切断は種数一である。
同定理の二重直線上の議論は $A(E+F)=3$ と標数零だけで一般切断が核空間に
閉埋込みされることを示すので、$p_a(J)=2$ に反する。

滑らかな二次曲面上の $A$ の型は和五、自乗八以上より $(1,4),(2,3)$。
前者は基点なしで核空間の一般切断は滑らかな有理曲線となる。
錐では $L=E+2F$, $A=aE+5F$ で $A^2\ge8$, $AE\ge0$ より $a=1,2$。
$a=1$ は基点なし、$AE=3$ である。逆系の $E$ 上の制限は基点なしで次数三なので、
一般切断は $E$ 上で異なる三点を切る。従って $J$ は錐頂点で正規化枝を三本以上持つ。
しかし前定理の証明の通り整Gorenstein種数二曲線には次数二の有限平坦な標準射があり、
正規化からの合成の各繊維の相異なる点は高々二つである。
残った両モデルで $A=2L+F$, $A^2=12$ である。

一般切断 $J_0=(a=0)$ は $S_0$ 上で整であり、核空間の $J$ に同型に写る。
滑らかな二次曲面では自明で、錐でも $J_0 E=1$ であるため $E$ と滑らかな一点で横断し、
頂点での像も滑らかである（最小解消の表示 $w(1,t,t^2)$ を使う）。
逆系は固定曲線を持たず、Bertini既約性を使った。
一般平面上では $J=G\to E_H\subset\mathbb P^2$ が射であるので、逆系の基点イデアルは
$J_0$ 上で主となる。実際一般点で入力座標を表す四切断の比は $J_0$ 上で射を表し、
各点で像の単元座標を選べば他の切断はその正則関数倍となる。
その生成元を $b$ に持ち上げると局所的に共通イデアルは $(a,b)$。
従って局所完全交叉であり、長さは一般二切断の交点数と基点外の八点の差 $12-8=4$ である。

階数一軌跡が有限なら定理\ref{thm:cofactor-exclude-nineteen}の三点積分の議論を適用できる。
依存する事項を明記する。長さ四の局所完全交叉を吹いた $\Gamma$ で $H=A-E_Z$ は自由、
$\Gamma$ は局所完全交叉である。解消上の基点重複度は $\sum m_i^2=4$ により
四単純点または一二重点である。同定理の交点数の評価により入力の正次元繊維の像は高々三点。
$\chi(\Gamma,L+tH)=4(t+1)^2$ と最大Cohen--Macaulay性から有限な開部分で
$M_D=f_*L$ である。従って二十印のうち少なくとも十七個が孤立した長さ三の繊維となる。
同定理の長さ四基点の平滑化と三点Hilbertスキームの計算はこの二つの曲面モデルと
局所完全交叉の基点に対するものであり、余因子曲線の整性を使わない。
一般四点への変形で三点繊維長は十六だから矛盾する。
最後に定理\ref{thm:rank-one-skew-pair-exclusion}を用いる。
\end{proof}
\path{computations/verify_sectional_nineteen_model.py} は各核曲面での数値候補、
四基点の長さおよび有限階数一の場合の排除に使う積分を検算する。
